% ============================================================
%  TERM PROJECT – PROPOSAL (template)
%  Databases · THGA Bochum
%
%  How to use:
%    1. Copy this folder into your own repository.
%    2. Replace every  <<< ... >>>  placeholder with your content.
%    3. Build with  `make`  (see the Makefile / README).
%    4. Send the resulting PDF to the lecturer and wait for approval.
% ============================================================
\documentclass[11pt,a4paper]{article}

\usepackage{thga-db}
\usepackage{caption}   % enables \captionof inside minipages

% ----- Fill in your metadata -------------------------------
\renewcommand{\thgaKurs}{Databases}
\renewcommand{\thgaDatum}{Summer Term 2026}
\newcommand{\authorName}{Abdelillah Taleb}
\newcommand{\projectTitle}{Hotel Booking Manager}

\begin{document}
\thispagestyle{empty}
\thgaTitel{Term Project – Proposal}{\projectTitle}

\begin{infobox}
\textbf{Author:} \authorName \\
\textbf{Module:} Introduction to Database Management Systems · THGA Bochum \\
\textbf{Submitted to:} Stephan Bökelmann · \href{mailto:sboekelmann@ep1.rub.de}{sboekelmann@ep1.rub.de} \\
\textbf{Target size:} about 6 pages
\end{infobox}

\begin{infobox}
\textbf{Every figure must be explained in the running text.} A figure never
stands on its own: refer to it by number and say what it shows -- e.g.
``Figure~\ref{fig:er} shows the ER model of \dots''.

Every figure below is \texttt{example-figure.jpg}, a dummy image shipped with
this template. Photograph or scan your own sketch, drop the file next to
\texttt{proposal.tex}, and point \texttt{\textbackslash includegraphics} at it:
\begin{quote}
\texttt{\textbackslash includegraphics[width=0.85\textbackslash linewidth]\{er-sketch.jpg\}}
\end{quote}
The \texttt{width} option scales the image to a fraction of the text width; the
file extension may be omitted. Keep the \texttt{\textbackslash caption} and the
\texttt{\textbackslash label} as they are.
\end{infobox}

% ============================================================
\section{Application domain and motivation}
% ============================================================

Hotel Booking Management System is a database-backed desktop
application designed for managing room reservations in a small hotel.
Instead of using spreadsheets or handwritten reservation books, all booking
information is stored in a relational PostgreSQL database.

The application is mainly used by hotel reception staff. It allows employees
to manage guests, rooms, bookings and additional hotel services such as
breakfast or parking. The system provides a structured and reliable way to
store reservation data while maintaining data integrity through a relational
database.

A relational database is well suited for this application because bookings
must always reference existing guests and rooms. Relationships between
entities can be enforced using primary keys and foreign keys, preventing
invalid or inconsistent data.

The first usage scenario is the creation of a new reservation. The receptionist
selects an available room for a given period, enters the guest information and
creates a booking.

The second usage scenario is managing an existing reservation. The receptionist
can view booking details, add optional services such as breakfast or parking
and update the booking status when the guest checks in or checks out.

% ============================================================
\section{Data model -- ER sketch}
% ============================================================

\begin{figure}[ht]
  \centering
  % Swap example-figure.jpg for your own file, e.g. er-sketch.jpg
  % (a photographed hand drawing is fine).
  \includegraphics[width=0.85\linewidth]{er-sketch.jpg}
  \caption{ER sketch of the Hotel Booking Manager data model, showing the
   central entity types, their key attributes, relationships and cardinalities.}
  \label{fig:er}
\end{figure}

Figure~\ref{fig:er} shows the ER model of the Hotel Booking Manager.
The central entity types are \texttt{guest}, \texttt{room},
\texttt{room\_category}, \texttt{booking}, \texttt{service} and
\texttt{booking\_service}.

Each guest is uniquely identified by a \texttt{guest\_id}. A guest can have
zero or many bookings, while every booking belongs to exactly one guest. Each
room is uniquely identified by a \texttt{room\_id}. A room can occur in many
bookings over time, while each booking refers to exactly one room.

Every room belongs to exactly one room category. A room category can describe
multiple rooms and stores shared information such as the category name,
maximum capacity and price per night. This avoids storing the same category
information repeatedly.

A booking can include several optional services, such as breakfast, parking or
late check-out. The same service can also be assigned to many different
bookings. Therefore, \texttt{booking} and \texttt{service} have an N:M
relationship. This relationship is resolved by the associative entity
\texttt{booking\_service}, which stores the selected service and its quantity
for a particular booking.

% ============================================================
\section{From model to schema}
% ============================================================

The ER model is transformed into the following relational schema. Each table
contains a primary key, and foreign keys are used to maintain the relationships
between the entities.

\begin{itemize}

\item \texttt{guest}(\underline{guest\_id}, first\_name, last\_name, email, phone)

\item \texttt{room\_category}(\underline{category\_id}, name, capacity, price\_per\_night)

\item \texttt{room}(\underline{room\_id}, room\_number, floor, active,
\emph{category\_id})

\item \texttt{booking}(\underline{booking\_id}, check\_in, check\_out,
status, \emph{guest\_id}, \emph{room\_id})

\item \texttt{service}(\underline{service\_id}, name, price)

\item \texttt{booking\_service}(\underline{booking\_id},
\underline{service\_id}, quantity)

\end{itemize}

\textbf{Normalisation.}

The database schema satisfies the Third Normal Form (3NF). Each table stores
information about exactly one entity, and all non-key attributes depend only on
the primary key. Repeated information, such as room categories and hotel
services, is stored only once and referenced through foreign keys, reducing
redundancy and improving data consistency.
% ============================================================
\section{API design}
% ============================================================

The backend of the Hotel Booking Manager is implemented using FastAPI and
provides REST endpoints for managing guests, rooms, bookings and services.
Read operations are available without authentication, while all endpoints that
modify data require a valid X-API-Key.

Some endpoints perform SQL JOIN operations to combine data from multiple
tables. An aggregation endpoint is included to calculate booking statistics.

\renewcommand{\arraystretch}{1.2}

\begin{tabularx}{\linewidth}{@{}p{2cm}p{5cm}X@{}}
\toprule
\textbf{Method} & \textbf{Path} & \textbf{Purpose} \\
\midrule

\texttt{GET}
&
\texttt{/rooms}
&
Returns all hotel rooms together with their room category (JOIN). \\

\texttt{GET}
&
\texttt{/bookings}
&
Returns all bookings including guest and room information (JOIN). \\

\texttt{GET}
&
\texttt{/statistics/bookings}
&
Returns the total number of bookings and occupied rooms (aggregation). \\

\texttt{POST}
&
\texttt{/bookings}
&
Creates a new booking (requires X-API-Key). \\

\texttt{PATCH}
&
\texttt{/bookings/\{id\}}
&
Updates an existing booking (requires X-API-Key). \\

\texttt{POST}
&
\texttt{/services}
&
Creates a new hotel service (requires X-API-Key). \\

\bottomrule
\end{tabularx}

% ============================================================
\section{Frontend prototype (hand sketches)}
% ============================================================

The frontend will be implemented as a desktop application using Python and
tkinter. It provides a simple and user-friendly interface for hotel reception
staff to manage guests, rooms, bookings and additional services.

When the application starts, the user enters the API URL and the X-API-Key in
a connection dialog. After a successful connection, the main screen displays
the available rooms and current bookings. Additional forms allow the user to
create new bookings, edit existing reservations and view booking details. The
frontend communicates with the FastAPI backend using REST requests.

% Each \includegraphics below still points at the dummy example-figure.jpg.
% Swap in your own file per screen, e.g. {fe-connect.jpg}, and keep the \captionof.
\begin{figure}[ht]
\centering
\begin{minipage}[t]{0.47\linewidth}\centering
  \includegraphics[width=\linewidth]{example-figure.jpg}
  \captionof{figure}{Startup connection dialog for entering the API URL and X-API-Key.}
  \label{fig:fe-connect}
\end{minipage}\hfill
\begin{minipage}[t]{0.47\linewidth}\centering
  \includegraphics[width=\linewidth]{example-figure.jpg}
  \captionof{figure}{Main screen displaying hotel rooms and current bookings using \texttt{GET /bookings}.}
  \label{fig:fe-main}
\end{minipage}

\vspace{1em}

\begin{minipage}[t]{0.47\linewidth}\centering
  \includegraphics[width=\linewidth]{example-figure.jpg}
  \captionof{figure}{Booking form for creating a new reservation using \texttt{POST /bookings}.}
  \label{fig:fe-detail}
\end{minipage}\hfill
\begin{minipage}[t]{0.47\linewidth}\centering
  \includegraphics[width=\linewidth]{example-figure.jpg}
  \captionof{figure}{Booking details screen for updating reservations using \texttt{PATCH /bookings/\{id\}}.}
  \label{fig:fe-second}
\end{minipage}
\end{figure}

Figures~\ref{fig:fe-connect}--\ref{fig:fe-second} illustrate the planned user
interface of the Hotel Booking Manager.

Figure~\ref{fig:fe-connect} shows the startup connection dialog where the user
enters the API URL and the X-API-Key before accessing the system.
Figure~\ref{fig:fe-main} shows the main screen, where reception staff can view
hotel rooms and current bookings retrieved from the backend using the
\texttt{GET /bookings} endpoint.

Figure~\ref{fig:fe-detail} presents the booking form, which is used to create
new reservations by sending data to the \texttt{POST /bookings} endpoint.
Finally, Figure~\ref{fig:fe-second} shows the booking details screen, where
existing reservations can be updated using the
\texttt{PATCH /bookings/\{id\}} endpoint.

% ============================================================
\section{Operation with Docker Compose}
% ============================================================

The application is deployed using Docker Compose. The system consists of two
main services: a PostgreSQL database and a FastAPI backend. Docker Compose
starts both containers and automatically connects them through a shared Docker
network.

The PostgreSQL container stores all application data, including guests, rooms,
bookings and hotel services. The FastAPI container provides the REST API and
communicates directly with the database.

Configuration values such as the database username, password, database name and
API settings are stored in a separate \texttt{.env} file. This keeps sensitive
information outside the application source code and allows different
configurations for development and deployment.

The backend is deployed on the lecture server as a Docker Compose application.
After uploading the project files, the services can be started with
\texttt{docker compose up -d}. The frontend connects to the backend by using
the configured API URL and a valid X-API-Key.

% ============================================================
\section{Scope, milestones and risks}
% ============================================================

\textbf{Scope.}

The project consists of six database tables and six REST API endpoints.
The database manages guests, rooms, room categories, bookings, hotel services
and the relationship between bookings and services. The API supports reading,
creating and updating booking information.

\textbf{Milestones and schedule.} <<< Fill in the week and the estimated hours;
the hours should add up to roughly 40. >>>

\renewcommand{\arraystretch}{1.2}

\begin{tabularx}{\linewidth}{@{}p{1.4cm}Xp{1.8cm}@{}}
\toprule
\textbf{Week} & \textbf{Milestone} & \textbf{Est. hours} \\
\midrule

1 &
Design the database schema, create the ER model and implement PostgreSQL tables.
&
8 \\

2--3 &
Implement the FastAPI backend, database queries and X-API-Key authentication.
&
12 \\

4--5 &
Develop the tkinter frontend and connect it to the REST API.
&
10 \\

6 &
Testing, documentation, Docker deployment and project video.
&
10 \\

\bottomrule
\end{tabularx}

\textbf{Risks.}

One possible risk is incorrect handling of overlapping room bookings, which
could result in the same room being reserved for multiple guests at the same
time. Another risk is communication problems between the frontend, backend and
database during deployment with Docker Compose.

% ============================================================
\section{Acceptance criteria and testing}
% ============================================================

The project is considered complete when the database, backend and frontend
work together correctly. The following acceptance criteria define the expected
functionality of the system.

\textbf{Acceptance criteria.}

\begin{itemize}[topsep=2pt]
    \item A new booking can only be created if the selected guest and room
    exist in the database.
    
    \item The system prevents overlapping bookings for the same room during the
    same period.
    
    \item The booking overview correctly displays booking, guest and room
    information using SQL JOIN operations.
    
    \item The booking statistics endpoint returns the correct total number of
    bookings.
    
    \item All endpoints that modify data require a valid X-API-Key. Requests
    without a valid key are rejected with an authentication error.
\end{itemize}

\textbf{Test plan.}

\begin{itemize}[topsep=2pt]

\item \textbf{Database:}
Verify all primary keys, foreign keys, NOT NULL constraints and UNIQUE
constraints by inserting both valid and invalid test data.

\item \textbf{API:}
Test every REST endpoint using FastAPI's TestClient. Verify successful
requests as well as invalid input, missing resources and unauthorized access.

\item \textbf{Business rule:}
Test that the application rejects overlapping bookings for the same room and
time period while allowing bookings for different rooms.

\item \textbf{Frontend:}
Verify that the desktop application can connect to the API, display booking
information correctly and create new bookings successfully.

\end{itemize}

\vspace{1em}
\begin{infobox}
\textbf{Effort budget:} the whole term project (system, documentation and video)
should fit into roughly \textbf{40\,h of work} over at most \textbf{2 months}.
Keep your scope within that budget.
\end{infobox}
